\documentclass[11pt,a4paper]{article}
\usepackage[italian]{babel}
\usepackage[margin=2.5cm]{geometry}
\usepackage{parskip}
\usepackage{amsmath, amssymb, amsthm}
\usepackage{booktabs}
\usepackage{array}
\usepackage{xcolor}
\usepackage{needspace}
\usepackage{enumitem}
\usepackage{listings}
\usepackage{tikz}
\usepackage[most]{tcolorbox}
\usepackage[hidelinks]{hyperref}
\usetikzlibrary{decorations.pathreplacing, shapes.geometric, arrows.meta, positioning}

% Niente vedove/orfane: i paragrafi non lasciano righe isolate a fine/inizio pagina
\widowpenalty=10000
\clubpenalty=10000

% Elenchi più compatti
\setlist{itemsep=2pt, parsep=0pt, topsep=4pt, partopsep=0pt}

% --- Riquadri con bordo nero, senza numero (si spezzano tra due pagine) ---
% Uso: \begin{definizione} ... \end{definizione}
%      \begin{teorema}[Nome] ... \end{teorema}  (il nome è facoltativo)
\tcbset{nota/.style={enhanced jigsaw, breakable, colback=white,
  colframe=black, boxrule=0.5pt, sharp corners}}
\NewTColorBox{definizione}{}{nota,
  before upper={\textbf{Definizione.}\ }}
\NewTColorBox{teorema}{o}{nota,
  before upper={\textbf{Teorema\IfValueT{#1}{ (#1)}.}\ }}
\NewTColorBox{proposizione}{o}{nota,
  before upper={\textbf{Proposizione\IfValueT{#1}{ (#1)}.}\ }}
\NewTColorBox{proprieta}{o}{nota,
  before upper={\textbf{Proprietà\IfValueT{#1}{ (#1)}.}\ }}
\NewTColorBox{assioma}{o}{nota, boxrule=1pt,
  before upper={\textbf{Assioma\IfValueT{#1}{ (#1)}.}\ }}

% --- Ambienti semplici (senza riquadro) ---
% Il titolo sta in cima al blocco, anche se il contenuto inizia con un riquadro o
% con codice; e non resta mai da solo in fondo a una pagina.
% Uso: \begin{esempio}[Titolo facoltativo] ... \end{esempio}
\newcommand{\newrunin}[2]{%
  \NewDocumentEnvironment{#1}{o}{%
    \par\Needspace{4\baselineskip}\addvspace{\medskipamount}\noindent
    \textbf{#2}\IfValueT{##1}{ (##1)}\textbf{.}\ \ignorespaces}%
  {\par\addvspace{\medskipamount}}}
\newrunin{esempio}{Esempio}
\newrunin{esercizio}{Esercizio}
\newrunin{osservazione}{Osservazione}

% V e F nelle tavole di verità
\newcommand{\V}{\textsf{V}}
\newcommand{\F}{\textsf{F}}

% --- Codice C ---
% Uso: \begin{codice} ... \end{codice}      codice C numerato
%      \begin{comando} ... \end{comando}    comandi da terminale
%      \begin{risultato} ... \end{risultato} output di un programma
%      \lstinline|printf("ciao");|           codice dentro al testo
\lstdefinestyle{c}{
  language=C,
  basicstyle=\ttfamily\small,
  keywordstyle=\color{blue}\bfseries,
  commentstyle=\color{gray}\itshape,
  stringstyle=\color{red!70!black},
  numbers=left,
  numberstyle=\tiny\color{gray},
  numbersep=8pt,
  columns=flexible,
  keepspaces=true,
  breaklines=true,
  showstringspaces=false,
  tabsize=4,
  morekeywords={include, define},
  % lettere accentate nei commenti
  literate={à}{{\`a}}1 {è}{{\`e}}1 {é}{{\'e}}1 {ì}{{\`i}}1
           {ò}{{\`o}}1 {ù}{{\`u}}1 {È}{{\`E}}1 {À}{{\`A}}1
}
\lstset{style=c}
\newtcblisting{codice}{listing only, nota, left=6mm,
  listing options={style=c}}
\newtcblisting{comando}{listing only, nota,
  listing options={basicstyle=\ttfamily\small, numbers=none, language={},
    columns=flexible, keepspaces=true, breaklines=true}}
\newtcblisting{risultato}{listing only, enhanced jigsaw, breakable,
  colback=black!5, colframe=black, boxrule=0.5pt, sharp corners,
  listing options={basicstyle=\ttfamily\small, numbers=none, language={},
    columns=flexible, keepspaces=true, breaklines=true}}

% --- Diagrammi di flusso (simboli standard) ---
\tikzset{
  terminale/.style = {ellipse, draw, minimum width=2.5cm, minimum height=0.9cm},
  io/.style        = {trapezium, trapezium left angle=70, trapezium right angle=110,
                      draw, minimum width=2.5cm, minimum height=0.9cm},
  processo/.style  = {rectangle, draw, minimum width=2.5cm, minimum height=0.9cm},
  decisione/.style = {diamond, draw, aspect=2, minimum width=2.5cm},
  freccia/.style   = {thick, -{Stealth}}
}

\title{Gruppo moltiplicativo di $\mathbb{Z}_n$ e funzione di Eulero}
\author{Marco Cerbella}
\date{Lezione 14 -- 3 novembre 2026}

\begin{document}
\maketitle

\section{Il gruppo moltiplicativo di $\mathbb{Z}_n$}

\begin{itemize}
  \item $(\mathbb{Z}_n, \cdot)$ non è mai un gruppo, perché $[0]$ non è invertibile.
  \item $(\mathbb{Z}_p\setminus\{[0]\}, \cdot)$ con $p$ primo è un gruppo (infatti $(\mathbb{Z}_p, +, \cdot)$ è un campo).
  \item $(\mathbb{Z}_n\setminus\{[0]\}, \cdot)$ in generale non è un gruppo.
\end{itemize}
Possiamo però definire
\[
  \mathbb{Z}_n^* = \left\{ [a]_n : \operatorname{MCD}(a, n) = 1 \right\} \subset \mathbb{Z}_n .
\]
$(\mathbb{Z}_n^*, \cdot)$ è un gruppo, detto \emph{gruppo moltiplicativo di $\mathbb{Z}_n$}.

\begin{esempio}
\begin{itemize}
  \item $\mathbb{Z}_p^* = \{[1], \dots, [p-1]\}$, con $p$ primo.
  \item $\mathbb{Z}_6^* = \{[1], [5]\}$.
  \item $\mathbb{Z}_9^* = \{[1], [2], [4], [5], [7], [8]\}$.
\end{itemize}
\end{esempio}

\section{La funzione di Eulero}

\begin{definizione}[Funzione di Eulero]
Sia $n \ge 1$. La \emph{funzione di Eulero} $\varphi$ è
\[
  \varphi(n) = |\mathbb{Z}_n^*| = \Big|\big\{ a \mid 1 \le a \le n,\ \operatorname{MCD}(a, n) = 1 \big\}\Big| .
\]
\end{definizione}

\begin{esempio}
$\varphi(6) = 2$, \qquad $\varphi(9) = 6$.
\end{esempio}

\begin{osservazione}
$\varphi(p) = p - 1$ se $p$ è primo, e $\varphi(p^2) = p^2 - p$.
\end{osservazione}

\begin{proposizione}
Sia $p$ primo e $n \ge 1$ intero. Allora
\[
  \varphi(p^n) = p^n - p^{n-1} = p^{n-1}(p - 1) .
\]
\end{proposizione}

\begin{proof}
Tra $1$ e $p^n$ ci sono $p^n$ numeri. Quelli non coprimi con $p^n$ hanno la forma $k\cdot p$, e $k$ può avere tutti i valori da $1$ fino a $p^{n-1}$: ci sono $p^{n-1}$ possibili numeri non coprimi con $p^n$. Quindi $\varphi(p^n) = p^n - p^{n-1}$.
\end{proof}

\begin{esempio}
$\varphi(125) = \varphi(5^3) = 5^3 - 5^2 = 125 - 25 = 100$.
\end{esempio}

\section{La biiezione $\mathbb{Z}_{nk} \to \mathbb{Z}_n \times \mathbb{Z}_k$}

Consideriamo la funzione
\[
  \psi : \mathbb{Z}_{nk} \to \mathbb{Z}_n \times \mathbb{Z}_k, \qquad [a]_{nk} \mapsto \big([a]_n, [a]_k\big).
\]

\begin{esempio}
\begin{itemize}
  \item $\psi_1 : \mathbb{Z}_{12} \to \mathbb{Z}_3 \times \mathbb{Z}_4$:
  \[
    \begin{gathered}
    [0]_{12} \mapsto ([0]_3, [0]_4), \quad [1]_{12} \mapsto ([1]_3, [1]_4), \quad [3]_{12} \mapsto ([0]_3, [3]_4),\\
    [6]_{12} \mapsto ([0]_3, [2]_4), \quad [7]_{12} \mapsto ([1]_3, [3]_4).
    \end{gathered}
  \]
  $\psi_1$ è biunivoca.
  \item $\psi_2 : \mathbb{Z}_{12} \to \mathbb{Z}_6 \times \mathbb{Z}_2$:
  \[
    \begin{gathered}
    [0]_{12} \mapsto ([0]_6, [0]_2), \quad [1]_{12} \mapsto ([1]_6, [1]_2), \quad [3]_{12} \mapsto ([3]_6, [1]_2),\\
    [6]_{12} \mapsto ([0]_6, [0]_2), \quad [7]_{12} \mapsto ([1]_6, [1]_2).
    \end{gathered}
  \]
  $\psi_2$ \emph{non è iniettiva} (ad esempio $[1]_{12}$ e $[7]_{12}$ hanno la stessa immagine) e \emph{non è suriettiva}: non esiste $[a]_{12}$ tale che $\psi_2([a]_{12}) = ([0]_6, [1]_2)$.
\end{itemize}
\end{esempio}

\begin{osservazione}
$[a]_{nk}$ è invertibile in $\mathbb{Z}_{nk}$ se e solo se $[a]_n$ è invertibile in $\mathbb{Z}_n$ e $[a]_k$ è invertibile in $\mathbb{Z}_k$.

\emph{Dimostrazione.} $\Rightarrow$ Se $[a]_{nk}$ è invertibile in $\mathbb{Z}_{nk}$, allora $\operatorname{MCD}(a, nk) = 1$, quindi $a$ non ha fattori comuni con $n$ e non ne ha con $k$, cioè $\operatorname{MCD}(a, n) = 1$ e $\operatorname{MCD}(a, k) = 1$: $[a]_n$ è invertibile in $\mathbb{Z}_n$ e $[a]_k$ è invertibile in $\mathbb{Z}_k$. $\Leftarrow$ Viceversa, se $\operatorname{MCD}(a, n) = 1$ e $\operatorname{MCD}(a, k) = 1$, allora $a$ non ha fattori comuni né con $n$ né con $k$, quindi $\operatorname{MCD}(a, nk) = 1$.
\end{osservazione}

\begin{proposizione}
Se $\operatorname{MCD}(n, k) = 1$, allora $\psi : \mathbb{Z}_{nk} \to \mathbb{Z}_n \times \mathbb{Z}_k$, $\psi([a]_{nk}) = ([a]_n, [a]_k)$, è \emph{biunivoca}.
\end{proposizione}

\begin{proof}
Si ha $|\mathbb{Z}_{nk}| = nk$ e $|\mathbb{Z}_n\times\mathbb{Z}_k| = |\mathbb{Z}_n|\cdot|\mathbb{Z}_k| = nk$. Proviamo che $\psi$ è iniettiva. Siano $[a]_{nk}, [b]_{nk}$ tali che $\psi([a]_{nk}) = \psi([b]_{nk})$, cioè $([a]_n, [a]_k) = ([b]_n, [b]_k)$, cioè
\[
  \begin{cases} [a]_n = [b]_n \\ [a]_k = [b]_k \end{cases}
  \quad\text{cioè}\quad
  \begin{cases} a \equiv b \pmod n \\ a \equiv b \pmod k \end{cases}
  \ \Rightarrow\ a \equiv b \pmod{\operatorname{mcm}(n, k)} .
\]
Ma $\operatorname{MCD}(n, k) = 1$, quindi $\operatorname{mcm}(n, k) = nk$, cioè $a \equiv b \pmod{nk}$, cioè $[a]_{nk} = [b]_{nk}$. Quindi $\psi$ è iniettiva; dominio e codominio hanno la stessa cardinalità finita, quindi $\psi$ è biunivoca.
\end{proof}

\begin{proposizione}[Moltiplicatività di $\varphi$]
Sia $\operatorname{MCD}(n, k) = 1$. Allora
\[
  \varphi(nk) = \varphi(n)\cdot\varphi(k) .
\]
\end{proposizione}

\begin{proof}
$\varphi(nk) = |\mathbb{Z}_{nk}^*| = |\mathbb{Z}_n^*\times\mathbb{Z}_k^*| = |\mathbb{Z}_n^*|\cdot|\mathbb{Z}_k^*| = \varphi(n)\cdot\varphi(k)$, dove si è usato che, per l'osservazione precedente, la restrizione $\psi : \mathbb{Z}_{nk}^* \to \mathbb{Z}_n^*\times\mathbb{Z}_k^*$ è biunivoca.
\end{proof}

\end{document}
